\subsection{Does every requested class change motion in the same way?}
Three secondary contrasts use the same perturbation for motorcycle, bicycle and taxi labels. The reference label in each comparison is private car. Each uses the requests originally belonging to its own two classes and retains only routes that both can legally use. All proposed motorcycle/car and taxi/car paths qualified. Only 48 of 96 bicycle/car paths qualified, illustrating the role of permissions even before motion synthesis. Since the path populations differ, the rows cannot be combined to rank all classes on one common network.

\fig{class-effects.pdf}{1}{Changing the class while preserving each road, start and initial noise. Left: change in travel. Right: change in the fraction of intervals below 0.1 m/s, in percentage points. All four contrasts are shown; bus/car is the primary replication. Each row uses its own jointly permitted paths.}{classeffects}
Relative to the car setting, the motorcycle setting increased travel by 146 metres (95\% interval 116 to 173) and reduced near-stationary time by 2.36 percentage points. The bicycle setting reduced travel by 577 metres (488 to 667 metres shorter) but also reduced near-stationary time by 5.15 points. A shorter trip can therefore arise through slower progression rather than more stopping. In the bicycle comparison 16 of the 2,304 outputs reached the path-length cap, all in the private-car condition. They remain in the totals, and endpoint clipping may contribute to the car stop fraction. The taxi setting reduced travel by 146 metres (92 to 203 metres shorter), with an uncertain change in near-stationary time ($-0.72$ points, interval $[-2.04,0.67]$). Every generated output in these comparisons passed the current road audit. The intervals are conditional on the eight simulation episodes and trained models, with no multiple-comparison adjustment.

The taxi result is instructive for the feedback question. A generative model can respond to a taxi label even when an auxiliary recognizer has difficulty identifying taxis from their completed traces. The two tasks have different observables. A change in generated distance alone says neither that an auditor should reliably recover the class nor that the distance distribution is calibrated to a particular taxi service.
